\chapter{Contemporary Issues 2021}

\section{The Sex Ratio at Birth: The Case Studies}

\begin{KeyConcept}
\begin{itemize}
\item \textbf{NFHS Round 5} key findings: overall sex ratio improved (1,020 women per 1,000 men), but the \textbf{sex ratio at birth (SRB) is 929} (up from 919 in NFHS-4) --- still below the WHO's natural 952, so \textbf{anthropogenic/social intervention} is at work. The \textbf{Total Fertility Rate (TFR) has fallen to $\sim$2.0} --- below replacement (2.1) --- which threatens to turn India's \textbf{demographic dividend into a demographic burden} (an ageing population; fewer young workers). Feticide/infanticide is now spreading even to tribal areas and Kerala.
\item \textbf{Reasons for the low SRB (sociological case studies)}:
\item \textbf{Ashish Bose} (1991 census): with \emph{more} development, urbanisation and westernisation there is \emph{more} feticide --- developed states (Punjab, Haryana, Western UP, Delhi/Noida, Himachal, Chandigarh) have high feticide because of \textbf{good transport + easy access to sex-determination technology + upper/middle-class position} (able to afford '\textbf{mobile doctors}' who come home); plus the \textbf{status of sons}, \textbf{globalisation/migration} (sons go abroad), and \textbf{westernisation $\rightarrow$ mindless consumerism $\rightarrow$ dowry}. He coined the \textbf{BIMARU districts} ('daughter-elimination male-aspiring rage for ultrasound' --- black holes in India's demographic transition).
\item \textbf{L. S. Vishwanath}: \textbf{Sanskritization} --- low castes emulate the upper castes (including dowry), so infanticide spreads downwards (E. G. Jenkinson, magistrate of Saharanpur in the 1830s, observed low castes like the Tuglas and Kolis adopting it). \textbf{More reservation $\rightarrow$ more Hinduization $\rightarrow$ more Sanskritization $\rightarrow$ more dowry/feticide}.
\item \textbf{Amartya Sen \& Jean Dreze}: the Chamars of UP --- sex ratio fell from 986 (1901) to 880 (1981), so SCs are now 'more like the higher castes'.
\item \textbf{Satish Agnihotri}: \textbf{feticide} (prenatal, done 'scientifically' by professional doctors) is treated as \textbf{'holier' than infanticide} (barbaric) --- so sex-determination technology \textbf{legitimises} daughter-elimination (doctors claim to be doing 'social service'), and the practice spreads to those who never did it before. Rural areas, unable to afford the technology, use infanticide.
\item \textbf{Mira Shiva} (an \textbf{interactionist} perspective): women themselves opt for feticide --- not because they are heartless but to \textbf{save their daughters} from a patriarchal society (eve-teasing, molestation, bride-burning, dowry death).
\end{itemize}
\end{KeyConcept}

\begin{QuickRevision}
\begin{itemize}
\item NFHS-5: SRB 929 (improvement but below WHO 952); TFR $\sim$2.0 (demographic dividend $\rightarrow$ burden); feticide spreading to tribal areas/Kerala.
\item Ashish Bose (development $\rightarrow$ more feticide; mobile doctors; BIMARU districts), Vishwanath (Sanskritization), Sen \& Dreze (Chamars 986$\rightarrow$880), Agnihotri (feticide 'holier' than infanticide), Mira Shiva (interactionist --- women saving daughters).
\end{itemize}
\end{QuickRevision}

\section{Missing Women and the Gender Gap}

\begin{KeyConcept}
\begin{itemize}
\item \textbf{'Missing women'} (Amartya Sen) refers to the suspiciously low ratio of women to men \textbf{at birth} in parts of the developing world --- caused by \textbf{social discrimination} (food, education, healthcare, sex-selective abortion), \emph{not} biology. (Emily Oster's counter-claim --- that Hepatitis B caused it biologically --- is rejected.)
\item \textbf{Siwan Anderson \& Ray}: 'missing women' is really \textbf{missing adult women} --- the main cause is \textbf{widowhood}: a widow in rural India is seen as a \textbf{witch}, has no access to property/resources, faces stigma (bad-luck omen), seclusion, change of diet/dress, discouragement from remarriage --- so she is vulnerable to \textbf{rape and suicide}.
\item The \textbf{Global Gender Gap Index 2021}: India's gap \textbf{widened to 62.5\%} --- because of women's inadequate representation in politics/technical/leadership roles, the declining female labour-force participation, poor healthcare, the female--male literacy gap, and income inequality. Its four dimensions: \textbf{economic participation, educational attainment, health \& survival, political empowerment}.
\end{itemize}
\end{KeyConcept}

\begin{QuickRevision}
\begin{itemize}
\item Missing women (Sen) = low sex ratio at birth, social not biological (vs Emily Oster's Hepatitis B); Siwan Anderson \& Ray --- missing adult women (widowhood $\rightarrow$ witch, no property, suicide).
\item Global Gender Gap 2021: 62.5\%; four dimensions (economic, educational, health, political).
\end{itemize}
\end{QuickRevision}

\section{Dominant Castes and the Demand for Backward Status}

\begin{KeyConcept}
\begin{itemize}
\item Dominant castes (Maratha, Jat, Patel, Gujar, Kapu) now demand \textbf{backward status}. The \textbf{Maratha} reservation was declared \textbf{unconstitutional} by the Supreme Court (a state cannot decide backwardness --- that is for the National Commission for Backward Classes); the \textbf{Rohini Commission} proposes \textbf{sub-categorising OBCs} (the creamy layer has captured the benefits).
\item \textbf{Reasons for the demand}: (1) \textbf{agrarian distress} --- dominant (peasant) castes want non-farm occupations, but the private sector has no jobs, so they want government-job quotas --- \textbf{Satish Deshpande} calls it 'an expression of \textbf{caste entitlement} rather than a claim to class disadvantage' (a means to urbanisation); (2) \textbf{classicisation of caste} (D. L. Sheth) --- caste loses its traditional functions and acquires new economic interest/political identity (cf. Parsons's structural differentiation); (3) a \textbf{crippled urban economy}; (4) \textbf{relative deprivation} --- resentment at reservation-driven SC/ST upward mobility (reservation is for them a means to \textbf{preserve social distance}, not fix injustice); (5) \textbf{status inconsistency} (Gerhard Lenski) --- socially forward but economically backward, hence Deshpande's '\textbf{backward forwards}'; (6) \textbf{exclusion/inclusion error} --- the creamy layer captures benefits; (7) \textbf{political support} --- their numerical strength (vote banks).
\end{itemize}
\end{KeyConcept}

\begin{QuickRevision}
\begin{itemize}
\item Maratha reservation unconstitutional; Rohini Commission (OBC sub-categorisation).
\item Reasons: agrarian distress (Deshpande --- caste entitlement), classicisation (D. L. Sheth), crippled urban economy, relative deprivation, status inconsistency (Lenski --- 'backward forwards'), exclusion/inclusion error, political support.
\end{itemize}
\end{QuickRevision}

\section{Gail Omvedt and the Backward Class Movement}

\begin{KeyConcept}
\begin{itemize}
\item \textbf{Gail Omvedt} (who died this year) wrote \textbf{'Dalits and the Democratic Revolution'} and \textbf{'Cultural Revolt: Non-Brahmin Movement in Maharashtra'}.
\item Her comparison of pre- and post-independence movements: the pre-independence movements (Satya Shodak Samaj, the Self-Respect Movement, SNDP, the Justice Party, the Adi Dharma movement) were \textbf{anti-caste = anti-systemic movements} --- they sought to transform the basic structure of the caste system, and thereby \textbf{introduced democracy into India}. Their common features: \textbf{rejection of Hindu nationalism} (they proposed \textbf{Dalit-Bahujan nationalism}), \textbf{rejection of the Aryan racial theory} (Brahmins were not the original inhabitants --- Dalits/Dravidians were), \textbf{rejection of Hinduism} (an 'imposed' religion --- 'I was born Hindu, but I will not die Hindu'), \textbf{rejection of the Congress} (a 'Brahmin-bourgeois/Shetji-Bhatji party'), and being \textbf{economic radicals} (against moneylenders and capitalists).
\item They were led by \textbf{sub-elites} (elites of the downward castes --- Phule, Periyar, Ambedkar, Narayana Guru) who competed with the traditional (Brahmin) elites; they were \textbf{peasant-based mass movements of the Bahujan}; and they were a \textbf{cultural revolt} --- a \textbf{counter-hegemony} (Gramsci) against the cultural hegemony of the upper castes.
\end{itemize}
\end{KeyConcept}

\begin{QuickRevision}
\begin{itemize}
\item Gail Omvedt: Dalits and Democratic Revolution; Cultural Revolt.
\item Pre-independence movements = anti-caste/anti-systemic (introduced democracy); reject Hindu nationalism, Aryan theory, Hinduism, Congress; economic radicals.
\item Sub-elites; peasant-based Bahujan movements; cultural revolt (Gramsci's counter-hegemony).
\end{itemize}
\end{QuickRevision}

\section{Surrogacy, MTP and ART}

\begin{KeyConcept}
\begin{itemize}
\item Three laws dealing with \textbf{women's reproductive health}: the \textbf{Medical Termination of Pregnancy (MTP) Act}, the \textbf{Surrogacy (Regulation) Bill}, and the \textbf{ART (Assisted Reproductive Technology) Bill}.
\item The \textbf{surrogacy bill} recognises only two beneficiaries: the \textbf{'intending couple'} (a married heterosexual couple, roughly 35--45, if infertile) and the \textbf{'intending woman'} (divorced or widowed --- \textbf{not} a single/unmarried woman; she must marry to gain parenthood). The surrogate carries the couple's combined sperm and egg, so she cannot call herself the mother; \textbf{commercial surrogacy is banned}, leaving only \textbf{altruistic surrogacy}.
\item Sociological critique: the bill is \textbf{exclusionary} --- it \textbf{excludes lesbians, gays, bisexuals and transgender people from parenthood} (parenthood is a status), and it \textbf{excludes single mothers} --- thereby \textbf{reproducing heteronormative standards} and \textbf{reproducing marriage as an institution} (marriage is treated as so sacred that one may have to sacrifice to gain parenthood). This is a form of \textbf{social exclusion}.
\end{itemize}
\end{KeyConcept}

\begin{QuickRevision}
\begin{itemize}
\item Three reproductive-health laws: MTP, Surrogacy, ART.
\item Surrogacy: intending couple + intending woman (divorced/widowed --- not single); commercial banned (only altruistic).
\item Critique: excludes gay/lesbian/bisexual/transgender and single mothers $\rightarrow$ reproduces heteronormativity and marriage; social exclusion.
\end{itemize}
\end{QuickRevision}

\section{Democratic Backsliding, Wage for Housework and Further Issues}

\begin{KeyConcept}
\begin{itemize}
\item \textbf{Democratic backsliding} = the legal/constitutional 'un-democratising' of democracy --- e.g. the abolition of \textbf{Article 370} was legal and constitutional, but it is argued to be \textbf{undemocratic} because it did not take the confidence of all stakeholders into account. (Scholars: \textbf{Ashutosh Varshney}, and Levitsky \& Ziblatt of 'How Democracies Die'.) \textbf{Democratic resilience} = the institutions that keep democracy from collapsing despite backsliding --- \textbf{civil society, media, judiciary, the masses, and social media}.
\item \textbf{Wage for housework}: does paying women for housework \textbf{liberate} them, or \textbf{reproduce the status quo} (the state legitimising housework as 'women's deal')? Link it to the \textbf{waves of feminism}.
\item Further exam-relevant issues to read: \textbf{POCSO Act} (the 'skin-to-skin contact' judgment, later rectified by the Supreme Court), \textbf{LGBT marriage rights}, the \textbf{PM Kaushal Vikas Yojana} (third phase, $\sim$40\% women workforce), and \textbf{Startup India} (its first review article in EPW).
\end{itemize}
\end{KeyConcept}

\begin{QuickRevision}
\begin{itemize}
\item Democratic backsliding (legal but undemocratic --- Art. 370) vs democratic resilience (civil society, media, judiciary, masses, social media); scholars Varshney, Levitsky \& Ziblatt.
\item Wage for housework (liberate vs reproduce status quo; waves of feminism).
\item Also read: POCSO (skin-to-skin), LGBT marriage rights, PM Kaushal Vikas Yojana (40\% women), Startup India (EPW review).
\end{itemize}
\end{QuickRevision}
